% id: 176-12
% book: 베이직쎈 대수 · 10 등비수열
% section: 개념 60 등비수열의 합
% passage: 다음을 만족시키는 등비수열 $\{a_n\}$의 첫째항과 공비를 구하시오.
% figure: none
% answer: 
% answer_source: 
% variant_level: 0 (원본 전사)
% 이 파일은 build.mjs 가 items.json 에서 만든다. 고칠 때는 items.json 을 고친다.
\smash{\raisebox{1.5mm}{\dmkichul{베이직쎈 대수 176쪽 12번}}}
$S_3=21$, $S_6=189$
\dmhint{$S_n$, $S_{2n}$을 알면 $\dfrac{S_{2n}}{S_n}$을 구해.\\ 첫째항을 $a$, 공비를 $r$ $(r\ne 1)$라 하면\\ $S_3=21$에서 \quad $\dfrac{\blank(r^3-1)}{r-1}=21$ \quad $\cdots\cdots$ ㉠\\ $S_6=189$에서 \quad $\dfrac{a(r^{\text{\scriptsize\setlength{\fboxsep}{1.5pt}\framebox[1.5em]{\rule{0pt}{1.6ex}}}}-1)}{r-1}=189$\\ \hspace*{1.5em}$\therefore\ \dfrac{a(r^{\text{\scriptsize\setlength{\fboxsep}{1.5pt}\framebox[1.5em]{\rule{0pt}{1.6ex}}}}-1)(r^3+1)}{r-1}=189$ \quad $\cdots\cdots$ ㉡\\ ㉡$\div$㉠을 하면 \quad $r^3+1=\blank$, \quad $r^3=\blank$ \quad $\therefore\ r=\blank$\\ $r=\blank$를 ㉠에 대입하여 풀면 \quad $a=\blank$}
